\subsection{Экспериментальный конвейер и метрики качества}
\label{sec:2-5-pipeline-metrics}

После описания источника данных, предобработки облаков точек, взвешенного ICP и модели движения необходимо зафиксировать, как эти компоненты объединяются в единый вычислительный конвейер и как затем оценивается качество получаемых оценок. В данной работе программная реализация разделяет два этапа: сначала алгоритм лидарной одометрии запускается на выбранном коротком участке проезда и сохраняет покадровые оценки линейных и угловых скоростей, а затем по этим данным выполняется количественный и визуальный анализ. Такой способ организации эксперимента важен не только с вычислительной точки зрения, но и методологически: он позволяет сравнивать конфигурации алгоритма в одинаковых условиях и анализировать устойчивость не только в среднем по набору данных, но и на уровне отдельных дорожных эпизодов.

\subsubsection{Экспериментальный конвейер}

В качестве базовой единицы эксперимента в работе рассматривается не весь длинный маршрут целиком, а отдельный короткий отрезок проезда. Основной сценарий построен на нарезке тестовых последовательностей Boreas-RT на 10-секундные отрезки. При частоте лидара, используемой в наборе данных, это соответствует интервалам по 100 кадров. Такой выбор удобен по нескольким причинам. Во-первых, он позволяет получить большое число однотипно обработанных эпизодов и тем самым сделать сравнение конфигураций алгоритма более статистически содержательным. Во-вторых, короткий отрезок лучше подходит для анализа устойчивости: если на некотором участке происходит деградация оценки движения, то она не размывается усреднением по длинному маршруту, а остаётся локализованной в конкретном эпизоде.

Эксперимент проводится на выбранной 10-секундной части проезда, для которой задаётся диапазон кадров и затем запускается алгоритм лидарной одометрии. В ходе обработки для каждого кадра сохраняются временная метка, эталонные линейные и угловые скорости, а также их оценки, полученные алгоритмом. В результате формируется временной ряд оценок движения, пригодный для последующего количественного и визуального анализа.

Отдельный интерес в работе представляет сценарий перекрытия поля зрения крупным динамическим объектом. Как было показано в разделе~\ref{sec:2-2-preprocessing}, в облако точек искусственно вводится эффект частичной потери наблюдаемости сцены за счёт моделирования грузовика, движущегося по соседней полосе. В составе экспериментального конвейера этот сценарий задаётся параметрически: грузовик в начальный момент располагается с продольным смещением \(-5\)~м относительно лидара, а затем обгоняет беспилотный автомобиль с относительной скоростью 1~м/с. Тем самым на протяжении выбранной части проезда область перекрытия поля зрения изменяется во времени, что позволяет исследовать устойчивость алгоритма не в статической, а в динамически деградирующей сцене. Важно, что такой сценарий воспроизводится одинаково для всех рассматриваемых эпизодов, поэтому сравнение различных конфигураций алгоритма остаётся корректным.

Конвейер реализован как последовательная покадровая обработка выбранной части проезда. Для каждого кадра извлекаются облако точек, временная метка и эталонные скорости, после чего облако передаётся в модуль лидарной одометрии. Затем текущая оценка скоростей сохраняется вместе с эталонным значением; при необходимости дополнительно фиксируются число итераций ICP и признак сходимости.

\subsubsection{Метрики качества и визуальный анализ}

Ключевой особенностью принятой в работе схемы оценки является то, что метрики качества сначала вычисляются для одного отдельного отрезка проезда, а уже затем при необходимости усредняются по множеству отрезков. Иными словами, исходной единицей анализа является один 10-секундный эпизод. Для него по всем кадрам рассматриваются ошибки компонент линейной и угловой скорости, определяемые как разность между оценкой алгоритма и эталонным значением. Анализ проводится по шести компонентам: \(v_x\), \(v_y\), \(v_z\), \(\omega_x\), \(\omega_y\) и \(\omega_z\). Такой выбор позволяет выявлять, какие именно компоненты оказываются наиболее чувствительными к ухудшению наблюдаемости сцены.

После вычисления покадровых ошибок для одного отрезка проезда по каждой компоненте скорости строится набор агрегированных статистик. В работе используются средняя ошибка, стандартное отклонение, среднеквадратичная ошибка, средняя абсолютная ошибка, 95-й процентиль абсолютной ошибки и корень из среднеквадратичной ошибки. Средняя ошибка позволяет судить о наличии систематического смещения оценки. Стандартное отклонение характеризует разброс ошибок относительно среднего уровня. Метрики MAE и RMSE описывают типичный масштаб ошибки, причём RMSE сильнее чувствительна к редким крупным отклонениям. Наконец, 95-й процентиль абсолютной ошибки особенно важен в рассматриваемой задаче, поскольку позволяет оценить поведение алгоритма в хвосте распределения ошибок, то есть в тех кадрах, где устойчивость оценки движения нарушается сильнее всего.

Принципиально важно, что все эти показатели сначала определяются именно на уровне одного отрезка проезда. Такой подход позволяет не терять информацию о локальных срывах устойчивости. Длинный маршрут может включать участки с существенно различающейся геометрией сцены: например, открытые участки дороги, туннели, зоны с бедной структурой окружения или эпизоды активного перекрытия поля зрения грузовиком. Если сразу усреднять ошибки по большому набору данных, то различия между такими режимами частично сглаживаются. Напротив, оценка по отдельным отрезкам проезда позволяет увидеть, в каких именно дорожных условиях алгоритм работает устойчиво, а в каких начинает деградировать. Уже после этого метрики, вычисленные по отдельным отрезкам проезда, могут усредняться по множеству эпизодов, что даёт итоговую сводную картину качества для выбранной конфигурации алгоритма.

Численные метрики в работе дополняются визуальным анализом. Во-первых, для отдельных отрезков проезда строятся графики временных рядов, на которых сравниваются эталонные и оценённые компоненты линейной и угловой скорости. Такой способ представления позволяет увидеть не только величину ошибки, но и её структуру во времени: например, запаздывание оценки, кратковременные выбросы или участки систематического смещения. Во-вторых, по ошибкам строятся гистограммы распределений, позволяющие оценить симметрию распределения, наличие тяжёлых хвостов и общую форму шума. Это особенно полезно в контексте обсуждения фильтра Калмана из раздела~\ref{sec:2-4-motion-model-kalman}, поскольку даёт эмпирическое представление о том, насколько разумно приближение шумов к нормальному распределению. Наконец, визуальный анализ позволяет выделять отдельные выбросные отрезки проезда и затем разбирать их содержательно, сопоставляя ухудшение качества с конкретной дорожной сценой или фазой обгона грузовика.

Таким образом, в работе используется двухуровневая схема оценки качества. На первом уровне для каждого 10-секундного отрезка проезда формируются покадровые ошибки и вычисляются агрегированные метрики по компонентам скоростей. На втором уровне эти показатели сопоставляются и усредняются по множеству отрезков проезда, что позволяет получить обобщённую характеристику устойчивости алгоритма. Такое построение экспериментального конвейера и системы метрик хорошо согласуется с целью работы: исследовать, насколько устойчивый алгоритм лидарной одометрии сохраняет качество оценки движения беспилотного автомобиля в условиях перекрытия поля зрения.
